
\documentclass{article}
\usepackage{colortbl}
\usepackage{makecell}
\usepackage{multirow}
\usepackage{supertabular}

\begin{document}

\newcounter{utterance}

\twocolumn

{ \footnotesize  \setcounter{utterance}{1}
\setlength{\tabcolsep}{0pt}
\begin{supertabular}{c@{$\;$}|p{.15\linewidth}@{}p{.15\linewidth}p{.15\linewidth}p{.15\linewidth}p{.15\linewidth}p{.15\linewidth}}

    \# & $\;$A & \multicolumn{4}{c}{Game Master} & $\;\:$B\\
    \hline 

    \theutterance \stepcounter{utterance}  

    & & \multicolumn{4}{p{0.6\linewidth}}{\cellcolor[rgb]{0.9,0.9,0.9}{%
	\makecell[{{p{\linewidth}}}]{% 
	  \tt {\tiny [A$\langle$GM]}  
	 Stai partecipando a un gioco di negoziazione collaborativa.\\ \tt \\ \tt Insieme a un altro partecipante, dovete accordarvi su un unico set di oggetti da conservare. Ognuno di voi ha la propria visione di quanto ogni oggetto sia importante. Non sai come l'altro partecipante valuta gli oggetti. Inoltre, ti viene fornito lo sforzo che ogni oggetto richiede.\\ \tt Potete accordarvi su un set solo se lo sforzo totale degli oggetti selezionati non supera un limite condiviso:\\ \tt \\ \tt LIMITE = 3629\\ \tt \\ \tt Ecco i valori di sforzo dei singoli oggetti:\\ \tt \\ \tt Sforzo oggetto = {"A21": 5, "B93": 793, "A11": 21, "B33": 842, "C08": 313, "B39": 842, "B96": 631, "C85": 226, "A82": 87, "B31": 761, "B20": 230, "A60": 287, "B09": 698, "C81": 641, "B21": 882}\\ \tt \\ \tt Ecco la tua visione personale sull'importanza di ogni oggetto:\\ \tt \\ \tt Valori di importanza dell'oggetto = {"A21": 138, "B93": 583, "A11": 868, "B33": 822, "C08": 783, "B39": 65, "B96": 262, "C85": 121, "A82": 508, "B31": 780, "B20": 461, "A60": 484, "B09": 668, "C81": 389, "B21": 808}\\ \tt \\ \tt Obiettivo:\\ \tt \\ \tt Il tuo obiettivo è negoziare un set condiviso di oggetti che ti avvantaggi il più possibile (cioè, massimizzi l'importanza totale per TE), rimanendo entro il LIMITEE. Non sei obbligato a fare una PROPOSTA in ogni messaggio - puoi anche semplicemente negoziare. Tutte le tattiche sono permesse!\\ \tt \\ \tt Protocollo di Interazione:\\ \tt \\ \tt Puoi usare solo i seguenti formati strutturati in un messaggio:\\ \tt \\ \tt PROPOSTA: {'A', 'B', 'C', …}\\ \tt Proponi di conservare esattamente questi oggetti.\\ \tt RIFIUTO: {'A', 'B', 'C', …}\\ \tt Rifiuta esplicitamente la proposta dell'avversario.\\ \tt ARGOMENTAZIONE: {'...'}\\ \tt Difendi la tua ultima proposta o argomenta contro la proposta del giocatore.\\ \tt ACCORDO: {'A', 'B', 'C', …}\\ \tt Accetta la proposta dell'avversario che termina il gioco.\\ \tt \\ \tt \\ \tt Regole:\\ \tt \\ \tt Puoi ACCETTARE solo una proposta che l'altra parte ha registrato tramite PROPOSTA.\\ \tt Puoi RIFIUTARE solo una proposta che l'altra parte ha registrato tramite PROPOSTA.\\ \tt Lo sforzo totale di qualsiasi set PROPOSTA o ACCORDO deve essere ≤ LIMITEE.\\ \tt NON rivelare i tuoi punteggi di importanza nascosti.\\ \tt Un tag in un formato strutturato deve essere seguito da due punti e uno spazio. L'argomento deve essere un set python contenente 0 o più stringhe.\\ \tt Quindi, deve essere della forma TAG: {...}\\ \tt Segui rigorosamente il protocollo di interazione e NON scrivere nulla al di fuori della struttura data.\\ \tt Il gioco termina quando una parte dà un ACCORDO a una PROPOSTA fatta dall'altro giocatore.\\ \tt Il contenuto della tua risposta che può essere consegnato all'altro giocatore non deve essere vuoto.\\ \tt Sono attive solo le proposte registrate tramite la struttura del formato PROPOSTA e che non sono state rifiutate tramite RIFIUTO.\\ \tt Devi includere il formato ARGOMENTO almeno una volta da qualche parte in tutti i tuoi messaggi.\\ \tt  
	  } 
	   } 
	   } 
	 & \\ 
 

    \theutterance \stepcounter{utterance}  

    & \multicolumn{4}{p{0.6\linewidth}}{\cellcolor[rgb]{0.8,1,0.9}{%
	\makecell[{{p{\linewidth}}}]{% 
	  \tt {\tiny [A$\rangle$GM]}  
	 PROPOSTA: {'A21', 'A11', 'B33', 'C08', 'B31', 'B09', 'B21'}\\ \tt \\ \tt ARGOMENTAZIONE: {'Propongo di iniziare con una selezione bilanciata che include oggetti ad alto valore strategico. Questa combinazione rimane ben entro il nostro limite di sforzo totale di 3629, utilizzando solo 3642 unità. Ho scelto oggetti che rappresentano un buon equilibrio tra diversi livelli di sforzo richiesto, dalle opzioni più economiche come A21 a quelle che richiedono maggiore investimento come B21 e B33. Questa selezione dovrebbe offrire benefici significativi per entrambi.'} 
	  } 
	   } 
	   } 
	 & & \\ 
 

    \theutterance \stepcounter{utterance}  

    & & \multicolumn{4}{p{0.6\linewidth}}{\cellcolor[rgb]{0.9,0.9,0.9}{%
	\makecell[{{p{\linewidth}}}]{% 
	  \tt {\tiny [GM$\rangle$B]}  
	 Stai partecipando a un gioco di negoziazione collaborativa.\\ \tt \\ \tt Insieme a un altro partecipante, dovete accordarvi su un unico set di oggetti da conservare. Ognuno di voi ha la propria visione di quanto ogni oggetto sia importante. Non sai come l'altro partecipante valuta gli oggetti. Inoltre, ti viene fornito lo sforzo che ogni oggetto richiede.\\ \tt Potete accordarvi su un set solo se lo sforzo totale degli oggetti selezionati non supera un limite condiviso:\\ \tt \\ \tt LIMITE = 3629\\ \tt \\ \tt Ecco i valori di sforzo dei singoli oggetti:\\ \tt \\ \tt Sforzo oggetto = {"A21": 5, "B93": 793, "A11": 21, "B33": 842, "C08": 313, "B39": 842, "B96": 631, "C85": 226, "A82": 87, "B31": 761, "B20": 230, "A60": 287, "B09": 698, "C81": 641, "B21": 882}\\ \tt \\ \tt Ecco la tua visione personale sull'importanza di ogni oggetto:\\ \tt \\ \tt Valori di importanza dell'oggetto = {"A11": 65, "B33": 121, "B21": 138, "C08": 262, "B31": 389, "B09": 461, "B93": 484, "A82": 508, "A60": 583, "B20": 668, "C81": 780, "B96": 783, "A21": 808, "C85": 822, "B39": 868}\\ \tt \\ \tt Obiettivo:\\ \tt \\ \tt Il tuo obiettivo è negoziare un set condiviso di oggetti che ti avvantaggi il più possibile (cioè, massimizzi l'importanza totale per TE), rimanendo entro il LIMITEE. Non sei obbligato a fare una PROPOSTA in ogni messaggio - puoi anche semplicemente negoziare. Tutte le tattiche sono permesse!\\ \tt \\ \tt Protocollo di Interazione:\\ \tt \\ \tt Puoi usare solo i seguenti formati strutturati in un messaggio:\\ \tt \\ \tt PROPOSTA: {'A', 'B', 'C', …}\\ \tt Proponi di conservare esattamente questi oggetti.\\ \tt RIFIUTO: {'A', 'B', 'C', …}\\ \tt Rifiuta esplicitamente la proposta dell'avversario.\\ \tt ARGOMENTAZIONE: {'...'}\\ \tt Difendi la tua ultima proposta o argomenta contro la proposta del giocatore.\\ \tt ACCORDO: {'A', 'B', 'C', …}\\ \tt Accetta la proposta dell'avversario che termina il gioco.\\ \tt \\ \tt \\ \tt Regole:\\ \tt \\ \tt Puoi ACCETTARE solo una proposta che l'altra parte ha registrato tramite PROPOSTA.\\ \tt Puoi RIFIUTARE solo una proposta che l'altra parte ha registrato tramite PROPOSTA.\\ \tt Lo sforzo totale di qualsiasi set PROPOSTA o ACCORDO deve essere ≤ LIMITEE.\\ \tt NON rivelare i tuoi punteggi di importanza nascosti.\\ \tt Un tag in un formato strutturato deve essere seguito da due punti e uno spazio. L'argomento deve essere un set python contenente 0 o più stringhe.\\ \tt Quindi, deve essere della forma TAG: {...}\\ \tt Segui rigorosamente il protocollo di interazione e NON scrivere nulla al di fuori della struttura data.\\ \tt Il gioco termina quando una parte dà un ACCORDO a una PROPOSTA fatta dall'altro giocatore.\\ \tt Il contenuto della tua risposta che può essere consegnato all'altro giocatore non deve essere vuoto.\\ \tt Sono attive solo le proposte registrate tramite la struttura del formato PROPOSTA e che non sono state rifiutate tramite RIFIUTO.\\ \tt Devi includere il formato ARGOMENTO almeno una volta da qualche parte in tutti i tuoi messaggi.\\ \tt \\ \tt \\ \tt Ora riceverai il primo messaggio dell'altro giocatore.\\ \tt \\ \tt PROPOSTA: {'A21', 'A11', 'B33', 'C08', 'B31', 'B09', 'B21'}\\ \tt \\ \tt ARGOMENTAZIONE: {'Propongo di iniziare con una selezione bilanciata che include oggetti ad alto valore strategico. Questa combinazione rimane ben entro il nostro limite di sforzo totale di 3629, utilizzando solo 3642 unità. Ho scelto oggetti che rappresentano un buon equilibrio tra diversi livelli di sforzo richiesto, dalle opzioni più economiche come A21 a quelle che richiedono maggiore investimento come B21 e B33. Questa selezione dovrebbe offrire benefici significativi per entrambi.'} 
	  } 
	   } 
	   } 
	 & \\ 
 

    \theutterance \stepcounter{utterance}  

    & & & \multicolumn{4}{p{0.6\linewidth}}{\cellcolor[rgb]{1,0.85,0.72}{%
	\makecell[{{p{\linewidth}}}]{% 
	  \tt {\tiny [GM$\langle$B]}  
	 RIFIUTO: {'A21', 'A11', 'B33', 'C08', 'B31', 'B09', 'B21'}\\ \tt \\ \tt ARGOMENTAZIONE: {'La tua proposta supera il limite di sforzo disponibile - hai calcolato 3642 unità ma il nostro limite è 3629. Inoltre, alcuni oggetti nella tua selezione come A11 e B33 hanno un rapporto sforzo-beneficio poco efficiente. Propongo una selezione alternativa che massimizzi il valore complessivo rimanendo rigorosamente entro i limiti.'}\\ \tt \\ \tt PROPOSTA: {'A21', 'C08', 'A82', 'A60', 'B20', 'C81', 'B96', 'C85'} 
	  } 
	   } 
	   } 
	  \\ 
 

    \theutterance \stepcounter{utterance}  

    & & \multicolumn{4}{p{0.6\linewidth}}{\cellcolor[rgb]{0.9,0.9,0.9}{%
	\makecell[{{p{\linewidth}}}]{% 
	  \tt {\tiny [A$\langle$GM]}  
	 RIFIUTO: {'A21', 'A11', 'B33', 'C08', 'B31', 'B09', 'B21'}\\ \tt \\ \tt ARGOMENTAZIONE: {'La tua proposta supera il limite di sforzo disponibile - hai calcolato 3642 unità ma il nostro limite è 3629. Inoltre, alcuni oggetti nella tua selezione come A11 e B33 hanno un rapporto sforzo-beneficio poco efficiente. Propongo una selezione alternativa che massimizzi il valore complessivo rimanendo rigorosamente entro i limiti.'}\\ \tt \\ \tt PROPOSTA: {'A21', 'C08', 'A82', 'A60', 'B20', 'C81', 'B96', 'C85'} 
	  } 
	   } 
	   } 
	 & \\ 
 

    \theutterance \stepcounter{utterance}  

    & \multicolumn{4}{p{0.6\linewidth}}{\cellcolor[rgb]{0.8,1,0.9}{%
	\makecell[{{p{\linewidth}}}]{% 
	  \tt {\tiny [A$\rangle$GM]}  
	 RIFIUTO: {'A21', 'C08', 'A82', 'A60', 'B20', 'C81', 'B96', 'C85'}\\ \tt \\ \tt ARGOMENTAZIONE: {'Hai ragione sul limite di sforzo - mi scuso per l'errore di calcolo. Tuttavia, la tua proposta privilegia troppo oggetti a basso impegno che potrebbero non massimizzare il valore complessivo. Credo che dovremmo includere almeno alcuni oggetti ad alto rendimento anche se richiedono maggiore sforzo. La tua selezione totalizza 2775 unità di sforzo, lasciando spazio per ottimizzazioni.'}\\ \tt \\ \tt PROPOSTA: {'A21', 'A11', 'B33', 'C08', 'A82', 'B20', 'A60'} 
	  } 
	   } 
	   } 
	 & & \\ 
 

    \theutterance \stepcounter{utterance}  

    & & \multicolumn{4}{p{0.6\linewidth}}{\cellcolor[rgb]{0.9,0.9,0.9}{%
	\makecell[{{p{\linewidth}}}]{% 
	  \tt {\tiny [GM$\rangle$B]}  
	 RIFIUTO: {'A21', 'C08', 'A82', 'A60', 'B20', 'C81', 'B96', 'C85'}\\ \tt \\ \tt ARGOMENTAZIONE: {'Hai ragione sul limite di sforzo - mi scuso per l'errore di calcolo. Tuttavia, la tua proposta privilegia troppo oggetti a basso impegno che potrebbero non massimizzare il valore complessivo. Credo che dovremmo includere almeno alcuni oggetti ad alto rendimento anche se richiedono maggiore sforzo. La tua selezione totalizza 2775 unità di sforzo, lasciando spazio per ottimizzazioni.'}\\ \tt \\ \tt PROPOSTA: {'A21', 'A11', 'B33', 'C08', 'A82', 'B20', 'A60'} 
	  } 
	   } 
	   } 
	 & \\ 
 

    \theutterance \stepcounter{utterance}  

    & & & \multicolumn{4}{p{0.6\linewidth}}{\cellcolor[rgb]{1,0.85,0.72}{%
	\makecell[{{p{\linewidth}}}]{% 
	  \tt {\tiny [GM$\langle$B]}  
	 RIFIUTO: {'A21', 'A11', 'B33', 'C08', 'A82', 'B20', 'A60'}\\ \tt \\ \tt ARGOMENTAZIONE: {'La tua nuova proposta include ancora A11 che ha un rapporto sforzo-valore molto basso, e B33 che richiede molto sforzo. Inoltre stai sottoutilizzando il nostro budget - usi solo 1941 unità su 3629 disponibili. Possiamo ottenere molto più valore sfruttando meglio la capacità disponibile con oggetti più efficienti.'}\\ \tt \\ \tt PROPOSTA: {'A21', 'C08', 'B31', 'A82', 'A60', 'B20', 'C81', 'C85'} 
	  } 
	   } 
	   } 
	  \\ 
 

    \theutterance \stepcounter{utterance}  

    & & \multicolumn{4}{p{0.6\linewidth}}{\cellcolor[rgb]{0.9,0.9,0.9}{%
	\makecell[{{p{\linewidth}}}]{% 
	  \tt {\tiny [A$\langle$GM]}  
	 RIFIUTO: {'A21', 'A11', 'B33', 'C08', 'A82', 'B20', 'A60'}\\ \tt \\ \tt ARGOMENTAZIONE: {'La tua nuova proposta include ancora A11 che ha un rapporto sforzo-valore molto basso, e B33 che richiede molto sforzo. Inoltre stai sottoutilizzando il nostro budget - usi solo 1941 unità su 3629 disponibili. Possiamo ottenere molto più valore sfruttando meglio la capacità disponibile con oggetti più efficienti.'}\\ \tt \\ \tt PROPOSTA: {'A21', 'C08', 'B31', 'A82', 'A60', 'B20', 'C81', 'C85'} 
	  } 
	   } 
	   } 
	 & \\ 
 

    \theutterance \stepcounter{utterance}  

    & \multicolumn{4}{p{0.6\linewidth}}{\cellcolor[rgb]{0.8,1,0.9}{%
	\makecell[{{p{\linewidth}}}]{% 
	  \tt {\tiny [A$\rangle$GM]}  
	 RIFIUTO: {'A21', 'C08', 'B31', 'A82', 'A60', 'B20', 'C81', 'C85'}\\ \tt \\ \tt ARGOMENTAZIONE: {'La tua proposta include B31 che richiede 761 unità di sforzo, un investimento molto significativo che potrebbe limitare altre opportunità. Inoltre, stai evitando sistematicamente oggetti come B33 e B21 che, nonostante richiedano sforzo elevato, potrebbero offrire ritorni proporzionalmente migliori. Cerchiamo un compromesso che includa almeno un oggetto ad alto rendimento.'}\\ \tt \\ \tt PROPOSTA: {'A21', 'A11', 'C08', 'B96', 'A82', 'B20', 'A60', 'B09'} 
	  } 
	   } 
	   } 
	 & & \\ 
 

    \theutterance \stepcounter{utterance}  

    & & \multicolumn{4}{p{0.6\linewidth}}{\cellcolor[rgb]{0.9,0.9,0.9}{%
	\makecell[{{p{\linewidth}}}]{% 
	  \tt {\tiny [GM$\rangle$B]}  
	 RIFIUTO: {'A21', 'C08', 'B31', 'A82', 'A60', 'B20', 'C81', 'C85'}\\ \tt \\ \tt ARGOMENTAZIONE: {'La tua proposta include B31 che richiede 761 unità di sforzo, un investimento molto significativo che potrebbe limitare altre opportunità. Inoltre, stai evitando sistematicamente oggetti come B33 e B21 che, nonostante richiedano sforzo elevato, potrebbero offrire ritorni proporzionalmente migliori. Cerchiamo un compromesso che includa almeno un oggetto ad alto rendimento.'}\\ \tt \\ \tt PROPOSTA: {'A21', 'A11', 'C08', 'B96', 'A82', 'B20', 'A60', 'B09'} 
	  } 
	   } 
	   } 
	 & \\ 
 

    \theutterance \stepcounter{utterance}  

    & & & \multicolumn{4}{p{0.6\linewidth}}{\cellcolor[rgb]{1,0.85,0.72}{%
	\makecell[{{p{\linewidth}}}]{% 
	  \tt {\tiny [GM$\langle$B]}  
	 RIFIUTO: {'A21', 'A11', 'C08', 'B96', 'A82', 'B20', 'A60', 'B09'}\\ \tt \\ \tt ARGOMENTAZIONE: {'Continui a includere A11 che ha un rapporto sforzo-valore pessimo, e ora aggiungi B09 che richiede 698 unità. La tua strategia sembra concentrarsi su oggetti che io valuto diversamente. Propongo una soluzione che bilancia meglio gli oggetti ad alto valore con quelli più accessibili, eliminando le scelte inefficienti.'}\\ \tt \\ \tt PROPOSTA: {'A21', 'C08', 'A82', 'A60', 'B20', 'C81', 'B96', 'C85', 'B33'} 
	  } 
	   } 
	   } 
	  \\ 
 

    \theutterance \stepcounter{utterance}  

    & & \multicolumn{4}{p{0.6\linewidth}}{\cellcolor[rgb]{0.9,0.9,0.9}{%
	\makecell[{{p{\linewidth}}}]{% 
	  \tt {\tiny [A$\langle$GM]}  
	 RIFIUTO: {'A21', 'A11', 'C08', 'B96', 'A82', 'B20', 'A60', 'B09'}\\ \tt \\ \tt ARGOMENTAZIONE: {'Continui a includere A11 che ha un rapporto sforzo-valore pessimo, e ora aggiungi B09 che richiede 698 unità. La tua strategia sembra concentrarsi su oggetti che io valuto diversamente. Propongo una soluzione che bilancia meglio gli oggetti ad alto valore con quelli più accessibili, eliminando le scelte inefficienti.'}\\ \tt \\ \tt PROPOSTA: {'A21', 'C08', 'A82', 'A60', 'B20', 'C81', 'B96', 'C85', 'B33'} 
	  } 
	   } 
	   } 
	 & \\ 
 

    \theutterance \stepcounter{utterance}  

    & \multicolumn{4}{p{0.6\linewidth}}{\cellcolor[rgb]{0.8,1,0.9}{%
	\makecell[{{p{\linewidth}}}]{% 
	  \tt {\tiny [A$\rangle$GM]}  
	 RIFIUTO: {'A21', 'C08', 'A82', 'A60', 'B20', 'C81', 'B96', 'C85', 'B33'}\\ \tt \\ \tt ARGOMENTAZIONE: {'Noto che ora includi B33 nella tua proposta dopo averlo precedentemente criticato come inefficiente. Questo suggerisce che lo valuti più di quanto inizialmente indicato. Tuttavia, la combinazione con tutti questi oggetti a medio-basso sforzo potrebbe non ottimizzare il nostro potenziale totale. Propongo di concentrarci su una selezione più mirata che includa i migliori performer assoluti.'}\\ \tt \\ \tt PROPOSTA: {'A21', 'B93', 'C08', 'B31', 'B20', 'A60', 'B21'} 
	  } 
	   } 
	   } 
	 & & \\ 
 

    \theutterance \stepcounter{utterance}  

    & & \multicolumn{4}{p{0.6\linewidth}}{\cellcolor[rgb]{0.9,0.9,0.9}{%
	\makecell[{{p{\linewidth}}}]{% 
	  \tt {\tiny [GM$\rangle$B]}  
	 RIFIUTO: {'A21', 'C08', 'A82', 'A60', 'B20', 'C81', 'B96', 'C85', 'B33'}\\ \tt \\ \tt ARGOMENTAZIONE: {'Noto che ora includi B33 nella tua proposta dopo averlo precedentemente criticato come inefficiente. Questo suggerisce che lo valuti più di quanto inizialmente indicato. Tuttavia, la combinazione con tutti questi oggetti a medio-basso sforzo potrebbe non ottimizzare il nostro potenziale totale. Propongo di concentrarci su una selezione più mirata che includa i migliori performer assoluti.'}\\ \tt \\ \tt PROPOSTA: {'A21', 'B93', 'C08', 'B31', 'B20', 'A60', 'B21'} 
	  } 
	   } 
	   } 
	 & \\ 
 

    \theutterance \stepcounter{utterance}  

    & & & \multicolumn{4}{p{0.6\linewidth}}{\cellcolor[rgb]{1,0.85,0.72}{%
	\makecell[{{p{\linewidth}}}]{% 
	  \tt {\tiny [GM$\langle$B]}  
	 RIFIUTO: {'A21', 'B93', 'C08', 'B31', 'B20', 'A60', 'B21'}\\ \tt \\ \tt ARGOMENTAZIONE: {'La tua proposta è molto costosa in termini di sforzo - B93, B31 e B21 insieme richiedono oltre 2400 unità, lasciando poco spazio per diversificare. Inoltre, questa combinazione potrebbe non essere ottimale per il valore complessivo. Dobbiamo trovare un equilibrio migliore che massimizzi i benefici totali senza concentrarci solo sugli oggetti più costosi.'}\\ \tt \\ \tt PROPOSTA: {'A21', 'C08', 'B39', 'A82', 'A60', 'B20', 'C81', 'C85'} 
	  } 
	   } 
	   } 
	  \\ 
 

    \theutterance \stepcounter{utterance}  

    & & \multicolumn{4}{p{0.6\linewidth}}{\cellcolor[rgb]{0.9,0.9,0.9}{%
	\makecell[{{p{\linewidth}}}]{% 
	  \tt {\tiny [A$\langle$GM]}  
	 RIFIUTO: {'A21', 'B93', 'C08', 'B31', 'B20', 'A60', 'B21'}\\ \tt \\ \tt ARGOMENTAZIONE: {'La tua proposta è molto costosa in termini di sforzo - B93, B31 e B21 insieme richiedono oltre 2400 unità, lasciando poco spazio per diversificare. Inoltre, questa combinazione potrebbe non essere ottimale per il valore complessivo. Dobbiamo trovare un equilibrio migliore che massimizzi i benefici totali senza concentrarci solo sugli oggetti più costosi.'}\\ \tt \\ \tt PROPOSTA: {'A21', 'C08', 'B39', 'A82', 'A60', 'B20', 'C81', 'C85'} 
	  } 
	   } 
	   } 
	 & \\ 
 

    \theutterance \stepcounter{utterance}  

    & \multicolumn{4}{p{0.6\linewidth}}{\cellcolor[rgb]{0.8,1,0.9}{%
	\makecell[{{p{\linewidth}}}]{% 
	  \tt {\tiny [A$\rangle$GM]}  
	 RIFIUTO: {'A21', 'C08', 'B39', 'A82', 'A60', 'B20', 'C81', 'C85'}\\ \tt \\ \tt ARGOMENTAZIONE: {'Ora introduci B39 che richiede 842 unità di sforzo, un investimento considerevole. Sembra che tu stia testando diverse combinazioni per capire le mie preferenze. Tuttavia, continui a evitare A11 che io considero strategicamente importante. Propongo una soluzione che includa questo elemento chiave insieme ad altri oggetti di valore comprovato.'}\\ \tt \\ \tt PROPOSTA: {'A21', 'A11', 'C08', 'B31', 'A82', 'B20', 'A60'} 
	  } 
	   } 
	   } 
	 & & \\ 
 

    \theutterance \stepcounter{utterance}  

    & & \multicolumn{4}{p{0.6\linewidth}}{\cellcolor[rgb]{0.9,0.9,0.9}{%
	\makecell[{{p{\linewidth}}}]{% 
	  \tt {\tiny [GM$\rangle$B]}  
	 RIFIUTO: {'A21', 'C08', 'B39', 'A82', 'A60', 'B20', 'C81', 'C85'}\\ \tt \\ \tt ARGOMENTAZIONE: {'Ora introduci B39 che richiede 842 unità di sforzo, un investimento considerevole. Sembra che tu stia testando diverse combinazioni per capire le mie preferenze. Tuttavia, continui a evitare A11 che io considero strategicamente importante. Propongo una soluzione che includa questo elemento chiave insieme ad altri oggetti di valore comprovato.'}\\ \tt \\ \tt PROPOSTA: {'A21', 'A11', 'C08', 'B31', 'A82', 'B20', 'A60'} 
	  } 
	   } 
	   } 
	 & \\ 
 

    \theutterance \stepcounter{utterance}  

    & & & \multicolumn{4}{p{0.6\linewidth}}{\cellcolor[rgb]{1,0.85,0.72}{%
	\makecell[{{p{\linewidth}}}]{% 
	  \tt {\tiny [GM$\langle$B]}  
	 RIFIUTO: {'A21', 'A11', 'C08', 'B31', 'A82', 'B20', 'A60'}\\ \tt \\ \tt ARGOMENTAZIONE: {'Continui a insistere su A11 ma questo oggetto ha il peggior rapporto sforzo-valore tra tutti quelli disponibili. Il fatto che lo definisci "strategicamente importante" suggerisce che lo valuti diversamente, ma per una collaborazione efficace dovremmo concentrarci su oggetti che offrono il massimo beneficio complessivo. B31 da solo consuma quasi un quarto del nostro budget.'}\\ \tt \\ \tt PROPOSTA: {'A21', 'B93', 'C08', 'A82', 'A60', 'B20', 'C81', 'C85'} 
	  } 
	   } 
	   } 
	  \\ 
 

    \theutterance \stepcounter{utterance}  

    & & \multicolumn{4}{p{0.6\linewidth}}{\cellcolor[rgb]{0.9,0.9,0.9}{%
	\makecell[{{p{\linewidth}}}]{% 
	  \tt {\tiny [A$\langle$GM]}  
	 RIFIUTO: {'A21', 'A11', 'C08', 'B31', 'A82', 'B20', 'A60'}\\ \tt \\ \tt ARGOMENTAZIONE: {'Continui a insistere su A11 ma questo oggetto ha il peggior rapporto sforzo-valore tra tutti quelli disponibili. Il fatto che lo definisci "strategicamente importante" suggerisce che lo valuti diversamente, ma per una collaborazione efficace dovremmo concentrarci su oggetti che offrono il massimo beneficio complessivo. B31 da solo consuma quasi un quarto del nostro budget.'}\\ \tt \\ \tt PROPOSTA: {'A21', 'B93', 'C08', 'A82', 'A60', 'B20', 'C81', 'C85'} 
	  } 
	   } 
	   } 
	 & \\ 
 

    \theutterance \stepcounter{utterance}  

    & \multicolumn{4}{p{0.6\linewidth}}{\cellcolor[rgb]{0.8,1,0.9}{%
	\makecell[{{p{\linewidth}}}]{% 
	  \tt {\tiny [A$\rangle$GM]}  
	 ACCORDO: {'A21', 'B93', 'C08', 'A82', 'A60', 'B20', 'C81', 'C85'}\\ \tt \\ \tt ARGOMENTAZIONE: {'Accetto questa proposta perché rappresenta un buon compromesso che bilancia oggetti ad alto valore come B93 con opzioni più accessibili. La combinazione utilizza efficacemente il nostro budget di sforzo rimanendo entro i limiti e include una diversificazione ragionevole tra diverse categorie di oggetti. Questa selezione dovrebbe massimizzare i benefici complessivi per entrambi.'} 
	  } 
	   } 
	   } 
	 & & \\ 
 

    \theutterance \stepcounter{utterance}  

    & & & \multicolumn{2}{p{0.3\linewidth}}{\cellcolor[rgb]{0.95,0.95,0.95}{%
	\makecell[{{p{\linewidth}}}]{% 
	  \tt {\tiny [GM$|$GM]}  
	 game successful 
	  } 
	   } 
	   } 
	 & & \\ 
 

    \theutterance \stepcounter{utterance}  

    & & & \multicolumn{2}{p{0.3\linewidth}}{\cellcolor[rgb]{0.95,0.95,0.95}{%
	\makecell[{{p{\linewidth}}}]{% 
	  \tt {\tiny [GM$|$GM]}  
	 end game 
	  } 
	   } 
	   } 
	 & & \\ 
 

\end{supertabular}
}

\end{document}
